\documentclass[11pt]{letter}

\usepackage[margin=0.5in]{geometry}
\usepackage{hyperref}

\begin{document}

\begin{letter}{}

\opening{Dear Hiring Team,}

I am writing to apply for the Python Engineer, New Grad position at Jobright.ai. With an M.Sc. in Computer Science from the University of Calgary and hands-on experience building Python software, Django REST APIs, data pipelines, and AI-enabled automation workflows, I am excited by the opportunity to contribute to backend services and production AI agents used by real job seekers.

What interests me most about Jobright.ai is the practical nature of the product. Job searching is a workflow I know closely, and I understand how much time can be lost to searching, matching, organizing, and tailoring applications. I am especially drawn to the chance to build reliable Python services, APIs, data pipelines, and AI-powered workflows that make that process faster and more guided for users.

My strongest language is Python, and I have used it across backend development, scientific software, machine-learning workflows, and automation projects. As a Back-End Developer at Asr Gooyesh Pardaz, I built Django REST APIs backed by PostgreSQL for NLP, text-to-speech, and speech-to-text services. I implemented authentication and verification flows, payment-related backend logic, and database-backed service-plan features. That experience gave me practical grounding in API design, relational data modeling, integration work, debugging, and writing maintainable backend code.

During my graduate research with BigGeo, Mitacs, and the University of Calgary, I built Python and C++ software for processing, storing, and retrieving large Earth observation datasets. The work required designing data structures, automating preprocessing pipelines, debugging incorrect outputs, measuring performance, and improving one C++ encoding workload from 233 seconds to 26 seconds through multithreading. Although the domain was geospatial research, the engineering habits are directly transferable to backend product work: make data usable, keep workflows reproducible, document assumptions, and improve reliability through careful iteration.

I have also been building an agentic job-search automation platform using Docker Compose, n8n, LLM-based workflows, Google Sheets and Drive integrations, SearXNG, and Telegram. This project is especially relevant to Jobright.ai because it combines job-search reasoning, workflow automation, external integrations, structured data, and user-facing decision support. It has made me more interested in AI systems that are useful in daily work, not just impressive as demos.

I would bring strong Python skills, backend/API experience, PostgreSQL familiarity, AI/ML foundations, and a careful communication style suited to remote collaboration. FastAPI, Flask, asynchronous production systems, and recommendation systems are areas where I would continue growing, but I am comfortable learning new tools quickly and contributing first where my Python, Django, data-pipeline, debugging, and automation experience are strongest.

Thank you for considering my application. I would welcome the opportunity to discuss how my background in backend engineering, AI workflows, and data-intensive software can contribute to Jobright.ai's engineering team.

\closing{Sincerely,

Amir Mirzai Golpayegani

\href{mailto:amirgolpaa24@gmail.com}{amirgolpaa24@gmail.com}}

\end{letter}

\end{document}
